\section{Introducción}\label{introducciuxf3n}

La gestión eficiente de inventarios y la planificación logística constituyen factores fundamentales para garantizar la continuidad operativa y la calidad del servicio en organizaciones que dependen del abastecimiento periódico de productos, es el caso de aquellas que operan redes de máquinas expendedoras, donde la disponibilidad oportuna de insumos impacta directamente en la satisfacción de los usuarios, la eficiencia de los procesos de distribución y el aprovechamiento de los recursos operativos \citep{grzybowska2020,mehmood2023}.

Sin embargo, muchas organizaciones continúan gestionando sus actividades de reabastecimiento mediante esquemas basados en calendarios fijos, inspecciones periódicas o intervenciones reactivas derivadas de incidencias reportadas por los clientes \citep{rosello2023}. Este tipo de estrategias presenta diversas limitaciones, cuando el reabastecimiento se realiza siguiendo rutas predeterminadas, sin considerar el comportamiento real del consumo, se pueden generar desplazamientos innecesarios hacia equipos que aún cuentan con niveles adecuados de inventario; mientras que algunas máquinas pueden agotar sus insumos antes de la siguiente visita, ocasionando interrupciones del servicio y afectando la experiencia del cliente \citep{aguas2025}. Estos eventos se traducen en un uso ineficiente de recursos logísticos, mayores costos operativos y una disminución de la calidad del servicio, desde la percepción de los clientes \citep{grzybowska2020}.

En los últimos años, la Inteligencia Artificial (IA) y el aprendizaje automático han demostrado su capacidad para apoyar la resolución de problemas relacionados a la gestión de inventarios, el pronóstico de demanda y la optimización logística \citep{puma2025}. El análisis de datos históricos permite identificar patrones de comportamiento, estimar necesidades futuras de abastecimiento y generar información útil para la toma de decisiones \citep{mehmood2023}. Asimismo, la combinación de modelos predictivos con algoritmos de optimización posibilita la planificación de rutas más eficientes, alineadas con las necesidades reales de operación \citep{grzybowska2020}.

En Bolivia el servicio de máquinas expendedoras de alimentos y bebidas es un modelo de negocio en crecimiento y todavía la oferta es limitada, se puede encontrar este servicio en aeropuertos, terminales de buses, hospitales, oficinas, edificios corporativos y espacios públicos de las principales ciudades del país: La Paz, Cochabamba y Santa Cruz \citep{datec2024}. Son pocas empresas que ofrecen este servicio como Amarket \citep{farmacorp2024}, muchas de ellas operan mediante la modalidad de comodato o venta directa para empresas e instituciones, es el caso de Smart Vending que además se encarga de toda la logística de las máquinas \citep{smartvending2022}, otras ofrecen estos servicios como una forma de promocionar sus equipos, pues en realidad son proveedores de estos equipos, como lo hacen Vending Bolivia y Nestlé Allegra \citep{vendingbolivia2026,nestle2026}.

En Sucre, la capital constitucional de Bolivia, este servicio es ofrecido por Cafetalino, una empresa que nace el año 2018, la cual ofrece autoservicio de bebidas calientes en base a café y que se constituye en pionera de este modelo de negocio en la ciudad. Inició sus actividades con unas pocas máquinas ubicadas en el aeropuerto de Alcantarí, hospitales, oficinas bancarias y centros educativos, desde entonces ha crecido permanentemente y a la fecha cuenta con una amplia red de máquinas expendedoras ubicadas en espacios públicos cerrados, que además de los ya mencionados incluyen mercados y centros comerciales, en los cuales por sus características, se tiene alguna restricción de horario; sin embargo, cuenta también con equipos ubicados en espacios abiertos, directamente al paso, hacia la acera, de manera que están disponibles las 24 horas del día.

En este contexto se desarrolla el presente proyecto, enfocado en la gestión de la red de máquinas dispensadoras de Cafetalino, empresa que de acuerdo con la información obtenida mediante entrevistas realizadas al gerente de la organización (Anexo~\ref{anexo-entrevista-al-gerente-de-cafetalino}) y a un operador responsable de las actividades de reabastecimiento (Anexo~\ref{anexo-entrevista-a-un-operador-de-cafetalino}), enfrenta situaciones recurrentes asociadas a visitas de reposición que resultan innecesarias, modificaciones extraordinarias de rutas y episodios de desabastecimiento que afectan la continuidad del servicio, condiciones que generan costos adicionales de operación y repercuten negativamente en la experiencia de los usuarios finales, debido a que maneja un sistema de reabastecimiento periódico programado de insumos y en ocasiones reactivo al recibir reportes de fallos en el servicio; además de no disponer de un indicador que le permita saber si una determinada máquina, necesita o no recarga de insumos.

Por lo expuesto, el problema central se puede enunciar de la siguiente manera:

El actual sistema de reabastecimiento periódico y de rutas fijas utilizado por Cafetalino genera un impacto negativo doble: por un lado, máquinas desabastecidas que prestan un servicio inadecuado y generan sentimientos negativos en sus clientes; y por otro, un gasto innecesario de tiempo y recursos logísticos al realizar recorridos fijos de revisión a equipos que no siempre requieren reposición de insumos.

Para atender esta problemática, se propone el desarrollo de un sistema inteligente de apoyo a la toma de decisiones basado en técnicas de aprendizaje automático y optimización logística. Esta solución utilizará los registros históricos de operación disponibles en Cafetalino para estimar las necesidades futuras de reposición de insumos de cada máquina y realizar la planificación dinámica de recorridos de abastecimiento; con ello, se busca mejorar la eficiencia operativa, reducir desplazamientos innecesarios y contribuir a la disminución de incidencias relacionadas con el desabastecimiento.

El alcance se limita al desarrollo de un Producto Mínimo Viable (MVP) enfocado en la predicción de los requerimientos de reabastecimiento y en la planificación de las rutas de recarga. La propuesta no contempla sensores físicos, dispositivos IoT ni monitoreo en tiempo real, sino que aprovecha la información histórica existente como fuente principal de conocimiento; esta delimitación responde a las necesidades identificadas en la organización y a las restricciones técnicas y económicas evidenciadas durante el análisis.

En consecuencia, el objetivo general del trabajo es mejorar la gestión del reabastecimiento de la red de máquinas dispensadoras de Cafetalino mediante un sistema basado en Inteligencia Artificial que prediga la demanda y permita la planificación dinámica de las rutas logísticas. La contribución esperada es un efecto observable sobre la operación: menos visitas innecesarias, menos episodios de desabastecimiento y una asignación más eficiente de los recursos logísticos. El detalle de este objetivo y de los objetivos específicos en que se despliega se presenta en el Capítulo 2.

\subsection{Justificación del Trabajo}\label{justificaciuxf3n-del-trabajo}

La relevancia de la propuesta radica en que aborda una problemática real mediante el aprovechamiento de datos de operación disponibles desde el 2018, transformando dicha información en un activo estratégico para apoyar la planificación de las actividades de abastecimiento. En el contexto organizacional, se espera contribuir a una gestión más eficiente de los recursos logísticos, proporcionando información que facilite la toma de decisiones de la empresa y del personal operativo.

Además, la implementación de este sistema tendrá un impacto directo en la modernización de la logística del mercado en Sucre y en Bolivia, un entorno donde no existen soluciones de gestión predictiva para redes de distribución y que se limita a métodos reactivos y de rutas fijas: promoverá la eficiencia en el transporte, evitará traslados redundantes y reducirá los costos de distribución. Desde una perspectiva tecnológica, el proyecto permitirá evaluar la viabilidad de integrar aprendizaje automático y optimización de rutas en un entorno de operación real, mostrando cómo la inteligencia artificial puede abordar problemas complejos de la cadena de suministro adaptándose a las particularidades geográficas locales y a la volatilidad del consumo sin depender de telemetría en tiempo real.

Adicionalmente, el proyecto posee relevancia académica al constituir un caso aplicado de utilización de técnicas de Inteligencia Artificial para resolver un problema de logística y abastecimiento en un contexto empresarial real. La investigación integra conceptos de análisis de datos, aprendizaje automático y optimización operativa dentro de una propuesta orientada a generar valor práctico y conocimiento aplicable en futuros proyectos de innovación tecnológica. Asimismo, los clientes se beneficiarán con una mayor disponibilidad de productos y una reducción de incidencias asociadas a la falta de insumos.

El proyecto tiene tres beneficiarios. Los clientes, que recibirán un servicio ininterrumpido y de mejor calidad, sin los episodios de desabastecimiento que hoy generan frustración y desconfianza. Los operadores de reabastecimiento, que dejarán de realizar recorridos innecesarios o redundantes y aprovecharán mejor su jornada. Y Cafetalino, que convertirá su registro histórico de recargas en un activo de alto valor y reducirá sus costos de operación al eliminar los gastos logísticos superfluos.

\subsection{Estructura en secciones}\label{estructura-en-secciones}

El presente documento se estructura de la siguiente manera: el Capítulo 1 presenta el contexto, la problemática y la justificación del proyecto; el Capítulo 2 establece el objetivo general y los objetivos específicos; el Capítulo 3 desarrolla la aplicación de Design Thinking, desde la investigación con los involucrados hasta la selección del prototipo; el Capítulo 4 describe la metodología de trabajo basada en Scrum y la filosofía LEAN; el Capítulo 5 detalla la implementación de la propuesta, con su planificación, presupuesto, construcción, despliegue y mantenimiento; el Capítulo 6 plantea el diseño experimental con el que se validará la solución en operación; y el Capítulo 7 contrasta los resultados con los objetivos planteados y señala las líneas de trabajo futuro. Los anexos recogen los instrumentos de investigación, los diagramas de los prototipos, las capturas de la gestión del proyecto y el backlog inicial.
